% Created: Jul, 25, 2026 15:11:56 by Thejaka Jayasinghe
\documentclass[dvipdfmx,11pt,a4paper,uplatex]{jsarticle}
\usepackage[dvipdfmx]{graphicx}
\usepackage{amsmath,amssymb}
\usepackage[margin=20truemm]{geometry}
\usepackage[hyphens]{url}
\usepackage{listings,jvlisting}
\lstset{
  basicstyle={\ttfamily},
  identifierstyle={\small},
  commentstyle={\smallitshape},
  keywordstyle={\small\bfseries},
  ndkeywordstyle={\small},
  stringstyle={\small\ttfamily},
  frame={tb},
  breaklines=true,
  columns=[l]{fullflexible},
  numbers=left,
  xrightmargin=0zw,
  xleftmargin=3zw,
  numberstyle={\scriptsize},
  stepnumber=1,
  numbersep=1zw,
  lineskip=-0.5ex
}

% For Tikz
% \usepackage[margin=0cm, bottom=0cm, footskip=0cm]{geometry}
% \usepackage{tikz}
% \usetikzlibrary{positioning,shapes,shapes.symbols,patterns,arrows}
% \usetikzlibrary{decorations.pathreplacing}
% \usetikzlibrary{calc}
% \pagestyle{empty}

% include packages
% \input{hyperref_jp}
% \input{cleveref_jp}
% \input{dec_math}
% \input{dec_text}
% \input{dec_smikawa}
% \input{siunitx_jp}

\title{}
\author{}

\begin{document}
\maketitle
\section{IGA 06: Order Reduction, Knot Removal and Least-Squares Projection}

This lecture introduces techniques for reducing the complexity of B-spline and Bézier representations while preserving the geometry as accurately as possible.

\subsection{Order Reduction}

Order reduction approximates a Bézier or B-spline curve of higher degree by a lower-degree curve.

Given a Bézier curve of degree $p_1$

\[
\mathbf{C}(u)=\sum_{i=0}^{p_1} B_i^{p_1}(u)\mathbf{P}_i,
\]

the objective is to approximate it by a curve of degree $p_2$, where

\[
p_2 < p_1.
\]

The reduced curve is

\[
\tilde{\mathbf{C}}(u)=
\sum_{i=0}^{p_2}
B_i^{p_2}(u)\mathbf{Q}_i,
\]

where the new control points $\mathbf{Q}_i$ minimize the approximation error.

Order reduction decreases the computational cost but introduces approximation errors because the original curve cannot generally be represented exactly with a lower polynomial degree.

\bigskip

\subsection{Knot Removal}

Knot removal reduces the number of knot spans by deleting unnecessary knots while maintaining the shape of the curve as closely as possible.

Removing knots

\begin{itemize}
\item decreases the number of elements,
\item reduces the number of control points,
\item simplifies the spline representation.
\end{itemize}

Since an exact representation is generally impossible after removing knots, the new control points are computed to minimize the geometric error.

Increasing the polynomial order before knot removal can improve the approximation quality.

\bigskip

\subsection{Least-Squares Projection}

Least-squares projection computes the best approximation of one spline space onto another by minimizing the squared error.

The approximation error is defined as

\[
E=
\int_{\Omega}
\left|
\mathbf{C}(u)-\tilde{\mathbf{C}}(u)
\right|^2du.
\]

The optimal control points satisfy the least-squares problem

\[
\min_{\mathbf{Q}} E.
\]

Compared with direct order reduction or knot removal, least-squares projection generally provides a more accurate approximation because it minimizes the global error over the entire curve.

\bigskip

\subsection{Key Points}

\begin{itemize}
\item \textbf{Order Reduction:} Lower polynomial degree with minimum approximation error.
\item \textbf{Knot Removal:} Fewer knot spans and control points while preserving the geometry.
\item \textbf{Order Elevation + Knot Removal:} Increasing the polynomial degree before removing knots often produces a better approximation.
\item \textbf{Least-Squares Projection:} Provides the optimal approximation by minimizing the global squared error.
\end{itemize}

\bigskip

\subsection{Practice Exercises}

The lecture demonstrates the following numerical procedures:

\begin{enumerate}
\item Approximate a cubic Bézier curve using a quadratic Bézier curve.
\item Approximate the same curve using a two-element quadratic B-spline.
\item Remove knots from a quadratic B-spline and compare the resulting geometry.
\item Perform order elevation followed by knot removal.
\item Compare direct approximation with least-squares projection.
\end{enumerate}
\end{document}
